\chapter{Discussion and Interpretation}\label{ch:discussion}
% GRUNDSATZ: Nur Ergebnisse/Messungen der Gruppe interpretieren, keine neuen Berechnungen.
% Analytische Werte: Aufgabe31LAB1.mlx (Datenblattwerte Faulhaber 3257G012CR)

\section{Comparison of Measurement, Model and Simulation}\label{sec:comparison}
The identified PT2 model reproduces the measured step response of the angular
velocity well (Figure~\ref{fig:identification}): both the rise and the
stationary value are matched, and the remaining deviations are mainly caused by
the noise of the measured signal. This noise results from the calculation of the
rotational speed from the position change of the incremental encoder between
two samples~\cite{laboranleitung}, so that the speed is only available in
discrete steps. The overdamped PT2 structure therefore describes the dynamic
behaviour of the drive sufficiently for the controller design, which is
consistent with the analytical model in Equation~\eqref{eq:gp2}. The same
approach was already successful for the preparatory data, where an additional
dead time was required to describe the delayed reaction of the system
(Appendix~\ref{app:prep}).

A comparison with the analytical model shows clear differences. Evaluating
Equation~\eqref{eq:gp2} with the datasheet values of the motor~\cite{faulhaber3257}
yields a static gain of \qty{52.36}{\radian\per\second\per\volt} and time
constants of \qty{4.5}{\milli\second} and \qty{0.18}{\milli\second}. The static
gain refers to the motor shaft, whereas the identified gain refers to the output
shaft; both are therefore only comparable when the transmission ratio $N$ is
taken into account, which brings them into the same order of magnitude. The time
constants, however, differ substantially: the dominant identified time constant
$\tau_2 = \qty{0.1030}{\second}$ is more than an order of magnitude larger than
the analytical one. The main reason is the assumption of an unloaded motor in the
analytical model. The real test bench additionally contains the gearbox, the
belt drive and the pointer, whose inertia and friction slow down the system
considerably. Furthermore, the motor driver and the real friction behaviour are
not part of the analytical model. The analytical model is thus suitable to
justify the model structure, while the parameters for the controller design have
to be identified experimentally.

% TODO: Vergleich Simulation (Figure 3.2) <-> Messung (Runs 1 und 2) ergänzen,
% sobald die Simulationsergebnisse aus PID_simulation.slx vorliegen.

\section{Controller Performance}\label{sec:performance}
The controller designed with the Ziegler--Nichols method stabilises the position
control loop, but the step response with the calculated parameters (run~2)
shows an overshoot (Figure~\ref{fig:run2}). This is expected, since the tuning
rules of Ziegler and Nichols aim at a fast response with a decay ratio of about
one quarter rather than at a response without overshoot~\cite{ziegler1942}. The
Ziegler--Nichols parameters are therefore a starting point that has to be
adjusted on the real system.

The most effective adjustment is the reduction of the integral gain by a factor
of 100 (run~1), after which no overshoot is visible (Figure~\ref{fig:run1}). This
can be explained by the structure of the plant: $G_\mathrm{p}(s)$ already
contains an integrator, so that a constant reference is reached without
remaining control error even without integral action. A large integral gain
mainly accumulates the control error during the rise and thus causes the
overshoot. A small integral gain is nevertheless useful, since it compensates
constant disturbances such as friction torques. This is confirmed by the
disturbance test (Figure~\ref{fig:disturbance}), in which the manually deflected
pointer is returned to the reference.

The variation of the gains shows the expected tendencies. A strongly reduced
proportional gain (runs~3 and~4) leads to a slow response that has not settled
at the end of the recording (Figures~\ref{fig:run3} and~\ref{fig:run4}). Since
the proportional part provides only a small control value in this case, the
response is presumably dominated by the derivative part directly after the step
and by the slowly growing integral part afterwards, which drives the position
beyond the reference. The derivative gain acts as damping: increasing it
(run~5) suppresses the overshoot, reducing it (run~6) causes an overshoot. An
increased proportional gain (runs~7 and~8) is possible, and in combination with
a larger derivative gain (run~8), a well-damped response is obtained.

In all runs, the control value reaches the saturation limit of \qty{12}{\volt}
directly after the reference step. During this phase, the system is driven with
the maximum voltage, so that the rise of the position is limited by the
actuator rather than by the controller gains. This explains why the initial rise
of the step responses with different gains is very similar
(Appendix~\ref{app:steps}), and it shows that the linear design does not fully
describe the real behaviour. The clamping anti-windup prevents the integrator
from accumulating during the saturation.

The tracking of the sinusoidal reference illustrates the trade-off in the choice
of the proportional gain (Table~\ref{tab:sine}). A higher gain reduces the
control error, but it also amplifies the measurement noise of the encoder
signal, which results in a strongly fluctuating control value at
\qty{800}{\volt\per\radian} and above (Figures~\ref{fig:sine800}
and~\ref{fig:sine1200}). Such a control value increases the load on the motor
and the motor driver; a proportional gain of about \qty{600}{\volt\per\radian}
is therefore a reasonable compromise for this reference signal.

The results are subject to several limitations. Each parameter set is measured
only once, so that the reproducibility is not verified, and the recording time
of \qty{1}{\second} is too short to assess the settling of runs~3 and~4. The
manual disturbance is not reproducible in magnitude and timing. In addition, the
plant model is identified from a single step amplitude of \qty{12}{\volt}, so
that nonlinear effects such as static friction are not captured.

\section{Conclusion}\label{sec:conclusion}
The DC motor test bench is successfully modelled, identified and controlled. The
analytical model confirms the PT2 structure of the drive, while the parameters
of the real system, including gearbox and belt drive, are obtained by a
least-squares fit of the measured step response. The Ziegler--Nichols method
yields a stable PID position controller, which, however, causes an overshoot on
the real system. Reducing the integral gain removes the overshoot, and a
combination of an increased proportional and derivative gain (run~8) gives a
fast and well-damped response. The experiments further show that the actuator
saturation dominates the behaviour after large reference steps and that the
achievable proportional gain is limited by the measurement noise. For further
investigations, repeated measurements, longer recording times and an
identification at several step amplitudes, including the friction behaviour,
are recommended.
